\section{Subset Sum Equal to Target}
\label{sec:dp-subset-sum-target}

\problemstatement{Given an array $\textit{arr}$ of $n$ non-negative
integers and a target $k$, determine whether \textbf{some} subset of
$\textit{arr}$ sums to exactly $k$.}

\begin{patternbox}
This is the Recursion chapter's ``does a
subsequence with sum $K$ exist'' recursion, revisited now that we know
to look for overlapping subproblems: the state
$(\textit{index}, \textit{remaining})$ is reached along many different
pick/not-pick paths, so caching by that pair of parameters converts the
$O(2^n)$ recursion into an $O(n \cdot k)$ DP.
\end{patternbox}

\subsection*{Deriving the recurrence}

Let $f(\textit{index}, \textit{target})$ be \texttt{true} if some subset
of $\textit{arr}[0..\textit{index}]$ sums to exactly $\textit{target}$.
At index $\textit{index}$, either this element is excluded
($f(\textit{index}-1, \textit{target})$ must hold), or it is included
(requiring $\textit{arr}[\textit{index}] \le \textit{target}$, and
$f(\textit{index}-1, \textit{target}-\textit{arr}[\textit{index}])$ must
hold):
\[
  f(\textit{index}, \textit{target}) = f(\textit{index}{-}1, \textit{target})
  \ \lor\ \big(\textit{arr}[\textit{index}] \le \textit{target} \ \land\
  f(\textit{index}{-}1, \textit{target}-\textit{arr}[\textit{index}])\big).
\]
Base case: $f(0, \textit{target}) = (\textit{target}=0) \lor
(\textit{target} = \textit{arr}[0])$.

\subsection*{Approach 1: Recursion}

\begin{lstlisting}
#include <bits/stdc++.h>
using namespace std;

bool subsetSumRecursive(int index, int target, vector<int>& arr) {
    if (target == 0) return true;
    if (index == 0) return arr[0] == target;

    bool notPick = subsetSumRecursive(index - 1, target, arr);
    bool pick = false;
    if (arr[index] <= target) {
        pick = subsetSumRecursive(index - 1, target - arr[index], arr);
    }
    return notPick || pick;
}

int main() {
    vector<int> arr = {1, 2, 3, 4};
    int n = arr.size();
    cout << (subsetSumRecursive(n - 1, 6, arr) ? "true" : "false") << endl;
    return 0;
}
\end{lstlisting}

\begin{complexitybox}[Approach 1 Complexity]
\textbf{Time.} $O(2^n)$.
\textbf{Space.} $O(n)$ recursion stack.
\end{complexitybox}

\subsection*{Approach 2: Memoization}

\begin{lstlisting}
#include <bits/stdc++.h>
using namespace std;

bool subsetSumMemo(int index, int target, vector<int>& arr,
                    vector<vector<int>>& memo) {
    if (target == 0) return true;
    if (index == 0) return arr[0] == target;
    if (memo[index][target] != -1) return memo[index][target];

    bool notPick = subsetSumMemo(index - 1, target, arr, memo);
    bool pick = false;
    if (arr[index] <= target) {
        pick = subsetSumMemo(index - 1, target - arr[index], arr, memo);
    }
    return memo[index][target] = (notPick || pick);
}

int main() {
    vector<int> arr = {1, 2, 3, 4};
    int n = arr.size(), k = 6;
    vector<vector<int>> memo(n, vector<int>(k + 1, -1));
    cout << (subsetSumMemo(n - 1, k, arr, memo) ? "true" : "false") << endl;
    return 0;
}
\end{lstlisting}

\begin{complexitybox}[Approach 2 Complexity]
\textbf{Time.} $O(n \cdot k)$.
\textbf{Space.} $O(n \cdot k)$ memo $+$ $O(n)$ recursion stack.
\end{complexitybox}

\subsection*{Approach 3: Tabulation}

\begin{lstlisting}
#include <bits/stdc++.h>
using namespace std;

bool subsetSumTabulation(vector<int>& arr, int k) {
    int n = arr.size();
    vector<vector<bool>> dp(n, vector<bool>(k + 1, false));

    for (int index = 0; index < n; index++) dp[index][0] = true;
    if (arr[0] <= k) dp[0][arr[0]] = true;

    for (int index = 1; index < n; index++) {
        for (int target = 1; target <= k; target++) {
            bool notPick = dp[index - 1][target];
            bool pick = false;
            if (arr[index] <= target) {
                pick = dp[index - 1][target - arr[index]];
            }
            dp[index][target] = notPick || pick;
        }
    }
    return dp[n - 1][k];
}

int main() {
    vector<int> arr = {1, 2, 3, 4};
    cout << (subsetSumTabulation(arr, 6) ? "true" : "false") << endl;
    return 0;
}
\end{lstlisting}

\subsection*{Dry run of the tabulation table}

\dryrun{Take $\textit{arr} = [1,2,3,4]$, $k=6$.}

\begin{center}
\begin{tabular}{c|ccccccc}
$\textit{index}\backslash\textit{target}$ & 0 & 1 & 2 & 3 & 4 & 5 & 6 \\ \hline
$0$ (val $1$) & T & T & F & F & F & F & F \\
$1$ (val $2$) & T & T & T & T & F & F & F \\
$2$ (val $3$) & T & T & T & T & T & T & T \\
$3$ (val $4$) & T & T & T & T & T & T & T
\end{tabular}
\end{center}

At $\textit{index}=2$ ($\textit{arr}[2]=3$), $\textit{target}=6$:
not-pick gives $\textit{dp}[1][6]=\text{F}$; pick gives
$\textit{dp}[1][3]=\text{T}$ (since $3\le6$). So
$\textit{dp}[2][6]=\text{T}$, already confirming a subset exists
($\{2,4\}$ sums to $6$, found even before considering index $3$). The
final answer, $\textit{dp}[3][6]=\text{T}$, confirms it.

\begin{complexitybox}[Approach 3 Complexity]
\textbf{Time.} $O(n \cdot k)$.
\textbf{Space.} $O(n \cdot k)$ table, no recursion stack.
\end{complexitybox}

\subsection*{Approach 4: Space Optimization}

\begin{lstlisting}
#include <bits/stdc++.h>
using namespace std;

bool subsetSumSpaceOptimized(vector<int>& arr, int k) {
    int n = arr.size();
    vector<bool> prev(k + 1, false);

    prev[0] = true;
    if (arr[0] <= k) prev[arr[0]] = true;

    for (int index = 1; index < n; index++) {
        vector<bool> current(k + 1, false);
        current[0] = true;
        for (int target = 1; target <= k; target++) {
            bool notPick = prev[target];
            bool pick = (arr[index] <= target) ? prev[target - arr[index]] : false;
            current[target] = notPick || pick;
        }
        prev = current;
    }
    return prev[k];
}

int main() {
    vector<int> arr = {1, 2, 3, 4};
    cout << (subsetSumSpaceOptimized(arr, 6) ? "true" : "false") << endl;
    return 0;
}
\end{lstlisting}

\begin{complexitybox}[Approach 4 Complexity]
\textbf{Time.} $O(n \cdot k)$.
\textbf{Space.} $O(k)$ for the rolling row, instead of $O(n \cdot k)$.
\end{complexitybox}

\begin{takeawaybox}
This boolean-OR recurrence, indexed by $(\textit{index},
\textit{target})$, is the foundation of every 0/1-knapsack-family
problem in this chapter -- Partition Equal Subset Sum, Count Subsets
with Sum K, and Count Partitions with a Given Difference are all this
\emph{exact} table, either queried at a derived target or with the OR
replaced by addition for counting. Internalize this one table fully
before moving on, since nothing new is derived from here through
Section~\ref{sec:dp-count-partitions-diff}.
\end{takeawaybox}
